\subsection{谱的一个分划}

用 \(\overline{Z}\) 表示子集 \(Z \cup \{-\infty, +\infty\}\) （\(\overline{R}\) 中的）。若 u 是一个核或余核有限维的线性映射，则称 u 的指标为由下式定义的元素 \(\text{ind}(u)\) （\(\overline{Z}\) 中的）

\[
\operatorname{ind} (u) = \dim \operatorname{Coker} (u) - \dim \operatorname{Ker} (u)
\]

（参见 III, p. 67 的第 6 目）。

定义 3. --- 设 E 是完备可赋范空间，u 是 E 的一个自同态。对每个 \(n \in \overline{Z}\) ，用 \(\mathrm{Sp}_{n}(u)\) 表示这样的复数 \(\lambda \in \mathrm{Sp}_{\mathrm{e}}(u)\) 所成的集：使 \(u - \lambda1_{E}\) 是严格态射，其核或余核有限维，且其指标为 \(n\) 。我们用 \(\mathrm{Sp}_{\omega}(u)\) 表示在 \(\mathrm{Sp}_{\mathrm{e}}(u)\) 中子集 \(\mathrm{Sp}_{n}(u)\) （\(n \in \overline{Z}\) ）之并的补集。

诸集 \(\mathrm{Sp}_{\mathrm{s}}(u)\) 、\(\mathrm{Sp}_{n}(u)\) （\(n \in \overline{Z}\) ）与 \(\mathrm{Sp}_{\omega}(u)\) 构成 u 的谱的一个分划。

E 的每个余核有限维的自同态都是严格的（III, p. 52, 引理 6）。E 的弗雷德霍姆自同态就是 E 的核与余核均有限维的自同态（III, p. 52, 命题 11）。集 \(\mathbf{C} - \mathrm{Sp}_{\mathrm{e}}(u)\) 由那些 \(\lambda \in \mathbf{C}\) （使 \(u - \lambda 1_{\mathrm{E}}\) 是 E 的里斯自同态的）组成，而这样的自同态是 E 的指标为 0 的弗雷德霍姆自同态。

于是，对 \(\lambda \in \mathbf{C}\) 及 \(n\in \mathbf{Z} - \{0\}\) ，我们有：

\(\lambda \in \mathbf{C} - \mathrm{Sp}(u) \iff u - \lambda 1_{\mathrm{E}}\) 是自同构；

\(\lambda \in \mathrm{Sp}_{\mathrm{s}}(u) \quad \Longleftrightarrow \quad u - \lambda 1_{\mathrm{E}} \text{ 是里斯自同态，但不是自同构；}\)

\(\lambda \in \mathrm{Sp}_{0}(u) \quad \Longleftrightarrow \quad u - \lambda 1_{\mathrm{E}}\) 是 E 的指标为 0 的弗雷德霍姆自同态但不是 E 的里斯自同态；

\(\lambda \in \mathrm{Sp}_n(u) \quad \Longleftrightarrow \quad u - \lambda 1_{\mathrm{E}}\) 是 E 的指标为 \(n\) 的弗雷德霍姆自同态（\(n \neq 0\)）；

\[
\lambda \in \mathrm{Sp} _ {- \infty} (u) \quad \Longleftrightarrow \quad u - \lambda 1 _ {\mathrm{E}} \text{ 是严格态射，其核无限维且其余核有限维；}
\]

\(\lambda \in \mathrm{Sp}_{+\infty}(u) \quad \Longleftrightarrow \quad u - \lambda 1_{\mathrm{E}}\) 是严格态射，其核有限维且其余核无限维；

\(\lambda \in \mathrm{Sp}_{\omega}(u) \quad \Longleftrightarrow \quad \text{要么 } u - \lambda 1_{\mathrm{E}} \text{ 不是严格态射，要么其核与余核都是无限维。}\)

注记. --- 用 \(\pi\) 表示从 \(\mathcal{L}(\mathrm{E})\) 到 E 的卡尔金代数 \(\mathcal{C}alk(\mathrm{E})\) 上的典范同态（参见 III, p. 75 的第 4 目）。\(\pi(u)\) 相对于代数 \(\mathcal{C}alk(\mathrm{E})\) 的谱有时称为 \(u\) 的稳定谱。其元素是那些复数 \(\lambda\) （使 \(u - \lambda 1_{\mathrm{E}}\) 不是 E 的弗雷德霍姆自同态的）（III, p. 73 的定理 3 的推论 1）。因此它等于 \(\mathrm{Sp}_{\omega}(u) \cup \mathrm{Sp}_{-\infty}(u) \cup \mathrm{Sp}_{+\infty}(u)\) 。

设 A 是 E 的与 u 交换的自同态所成的集。\(\pi(\mathrm{A})\) 在 \(\mathcal{C}alk(\mathrm{E})\) 中的闭包 B 是一个完备可赋范代数，且 \(\pi(u)\) 相对于 B 的谱就是 u 的本质谱 \(\mathrm{Sp}_{\mathrm{e}}(u)\) （III, p. 77, 命题 3）。把 I, p. 28 的命题 6 的推论应用于 \(\mathcal{C}alk(\mathrm{E})\) 的子代数 B，便知集 \(\mathrm{Sp}_{\mathrm{e}}(u)\) 是 \(\mathrm{Sp}_{\omega}(u) \cup \mathrm{Sp}_{-\infty}(u) \cup \mathrm{Sp}_{+\infty}(u)\) 与其补集的某些有界连通分支的并。

定理 1. --- 设 E 是完备可赋范空间，u 是 E 的一个自同态。

\begin{enumerate}
\def\labelenumi{\alph{enumi})}
\item
  集 \(\mathrm{Sp}_{\omega}(u)\) 是紧的。若 E 无限维，则它非空；
\item
  设 \(n \in \overline{Z}\) 。集 \(\mathrm{Sp}_{n}(u)\) 是 \(\mathbf{C} - \mathrm{Sp}_{\omega}(u)\) 的一族有界连通分支的并。它在 C 中是开的，且它在 C 中的边界包含于 \(\mathrm{Sp}_{\omega}(u)\) 。
\end{enumerate}

集 \(\mathbf{C} - \mathrm{Sp}_{\omega}(u)\) 由那些复数 \(\lambda \in \mathbf{C}\) （使 \(u - \lambda 1_{\mathrm{E}}\) 为核或余核有限维的严格态射的）组成。由 III, p. 67 的命题 11 与 III, p. 70 的命题 13，它是开的。于是集 \(\mathrm{Sp}_{\omega}(u)\) 是闭的。由于它有界，它是紧的。

我们来证明 \(b\) 。设 \(n \in \overline{\mathbf{Z}}\) 。集 \(\mathrm{Sp}_n(u)\) 包含于 \(\mathbf{C} - \mathrm{Sp}_{\omega}(u)\) 。设 U 是 \(\mathbf{C} - \mathrm{Sp}_{\omega}(u)\) 的一个与 \(\mathrm{Sp}_n(u)\) 相交的连通分支。由于映射 \(\lambda \mapsto \operatorname{ind}(u - \lambda 1_{\mathrm{E}})\) （从 \(\mathbf{C} - \mathrm{Sp}_{\omega}(u)\) 到 \(\overline{\mathbf{Z}}\) 中）是局部常数的（III, p. 68 的命题 12 的推论 1 及 III, p. 70 的命题 13 的推论 1），\(u - \lambda 1_{\mathrm{E}}\) 的指标等于 \(n\) （对一切 \(\lambda \in U\) ）。若 \(n \neq 0\) ，这就蕴含 U 包含于 \(\mathrm{Sp}_n(u)\) 。若 \(n = 0\) ，注意到集 U 是 \(\mathbf{C} - (\mathrm{Sp}_{\omega}(u) \cup \mathrm{Sp}_{-\infty}(u) \cup \mathrm{Sp}_{+\infty}(u))\) 的一个连通分支。由于 U 与 \(\mathrm{Sp}_0(u)\) 相交，从而与 \(\mathrm{Sp}_{\mathrm{e}}(u)\) 相交，由注记 2 得出集 U 包含于 \(\mathrm{Sp}_{\mathrm{e}}(u)\) ，从而包含于 \(\mathrm{Sp}_0(u)\) 。在所有情形，我们都得出结论：\(\mathrm{Sp}_n(u)\) 是 \(\mathbf{C} - \mathrm{Sp}_{\omega}(u)\) 的与 \(\mathrm{Sp}_n(u)\) 相交的那些连通分支的并。由于集 \(\mathrm{Sp}(u)\) 有界，这些连通分支必然有界。于是 \(\mathrm{Sp}_n(u)\) 在 \(\mathbf{C}\) 中是开的，且其边界包含于 \(\mathrm{Sp}_{\omega}(u)\) 。这就证明了 \(b\) 。

最后假设集 \(\mathrm{Sp}_{\omega}(u)\) 是空的。由 b)，诸集 \(\mathrm{Sp}_{n}(u)\) （\(n \in \overline{Z}\) ）此时也都是空的。于是我们有 \(\mathrm{Sp}(u) = \mathrm{Sp}_{\mathrm{s}}(u)\) 。因而 u 的谱是离散且紧的，故有限；又由于其所有点都有有限谱重数，向量空间 E 是有限维的（III, p. 83, 命题 2）。这就完成了 a) 的证明。

推论. --- a) 设 \(\Omega\) 是 \(\mathbf{C} - \mathrm{Sp}_{\omega}(u)\) 的无界连通分支。我们有 \(\Omega \cap \mathrm{Sp}(u) \subset \mathrm{Sp}_{\mathrm{s}}(u)\) ；

\begin{enumerate}
\def\labelenumi{\alph{enumi})}
\setcounter{enumi}{1}
\tightlist
\item
  每个附着于 \(\mathrm{Sp}_{\mathrm{s}}(u)\) 但不属于 \(\mathrm{Sp}_{\mathrm{s}}(u)\) 的点都属于 \(\mathrm{Sp}_{\omega}(u)\) 。
\end{enumerate}

断言 \(a\) 是定理 1 的断言 \(b\) 的直接推论。设 \(\lambda\) 是附着于 \(\mathrm{Sp}_{\mathrm{s}}(u)\) 但不属于 \(\mathrm{Sp}_{\mathrm{s}}(u)\) 的一点。它属于 \(u\) 的谱，因为后者是闭的。它不属于任何集 \(\mathrm{Sp}_{n}(u)\) （\(n \in \overline{\mathbf{Z}}\) ），因为这些集是开的（同前）且与 \(\mathrm{Sp}_{\mathrm{s}}(u)\) 不相交。于是我们有 \(\lambda \in \mathrm{Sp}_{\omega}(u)\) ，由此得 \(b\) 。

命题 3. --- 设 E 和 F 是完备可赋范空间，\(u\colon E\to F\) 与 \(v\colon F\to E\) 是连续线性映射。

\begin{enumerate}
\def\labelenumi{\alph{enumi})}
\item
  在 \(\mathbf{C} - \{0\}\) 上，集 \(\mathrm{Sp}(v \circ u)\) 与 \(\mathrm{Sp}(u \circ v)\) （相应地，\(\mathrm{Sp}_{\mathrm{s}}(v \circ u)\) 与 \(\mathrm{Sp}_{\mathrm{s}}(u \circ v)\) ，相应地，\(\mathrm{Sp}_{n}(v \circ u)\) 与 \(\mathrm{Sp}_{n}(u \circ v)\) （\(n \in \overline{\mathbf{Z}}\) ），相应地，\(\mathrm{Sp}_{\omega}(v \circ u)\) 与 \(\mathrm{Sp}_{\omega}(u \circ v)\) ）的迹相等；
\item
  设 \(\lambda\) 是 \(\mathrm{Sp}_{\mathrm{s}}(v \circ u)\) 中异于 0 的一个元素。\(\lambda\) 对 \(v \circ u\) 与对 \(u \circ v\) 的谱重数相等。
\end{enumerate}

设 \(\mu\) 是一个非零复数，且 \(n \in \overline{\mathbf{Z}}\) 。为使 \(\mu 1_{\mathrm{E}} - v \circ u\) 是一个自同构（相应地，一个里斯自同态，相应地，一个核或余核有限维且指标为 \(n\) 的严格态射），必须且只需 \(\mu 1_{\mathrm{F}} - u \circ v\) 也是（III, p. 49, 命题 10）。断言 \(a\) 于是由定义得出。

设 \(\lambda\) 是 \(\mathrm{Sp}_{\mathrm{s}}(v \circ u)\) 中异于 0 的一点。我们有

\[
\dim \operatorname{Ker} ((\lambda 1 _ {\mathrm{E}} - v \circ u) ^ {n}) = \dim \operatorname{Ker} ((\lambda 1 _ {\mathrm{F}} - u \circ v) ^ {n})
\]

（对一切 \(n \geqslant 0\) ；同前），故 \(\lambda\) 对 \(v \circ u\) 与 \(u \circ v\) 的谱重数相等。

